\scanpage{642}
\mspage{30}

stable par cochangement de base, $TC$\deleted{\unsure{f}} l'étant.) Pour
vérifier que l'on a un triple de Quillen
très spécial, il \deleted{suffit de}\inserted{reste, en vertu de 2.8, à} vérifier que

\[
 TF[I]\subset W(I)
\]

Ce qui \deleted{est} s'écrit

\[
 v^*(TF[I])\subset W(I_0),
\]

où

\[
 v\colon I_0\unsure{-|\longrightarrow}I
\]

est défini comme dans (1.12) (p. \unsure{10|20}). En fait,
on aura même

\[
 v^*(TF[I])=v^*v_*(TF(I_0))\subset
 \underset{\substack{\|\\ F(I_0)\cap W(I_0)}}{TF(I_0)},
\]

et la démonstration est duale de celle de

\[
 v^*v_!(TC(I_0))\unsure{\subset|\in} TC(I_0),
\]

(cf (2.3.\unsure{7|2}) p. 21), elle ne dépend pas des
questions de factorisation dans $M(I)$, et provient
ici de la stabilité de $TF$ par petits
produits. Cela prouve $1^\circ$.

\deleted{Une fois}
La démonstration de $2^\circ$) est maintenant
\inserted{la duale de} celle de (2.3. $2^\circ$). Comme cela nous fait
\deleted{\unsure{prouver}}\inserted{répéter} deux fois la $=$ démonstration, cela
montre qu'on n'a pas ``monté la \unsure{suite|sauce|boîte}'' \unsure{tout}
à fait \unsure{avec} assez de soin. Donc à titre
de \unsure{rattrappage|rattrapage}, j'ajoute un
% archive sheet number: 641
